% Auto-generated by Code2Latex
% Requires: longtable, booktabs packages

\begin{longtable}{p{\textwidth}}
\caption{Tools for single (Level 1)} \\
\endfirsthead

\multicolumn{1}{c}{\tablename\ \thetable{} -- continued from previous page} \\
\endhead

\multicolumn{1}{r}{Continued on next page} \\
\endfoot

\endlastfoot

\textbf{\texttt{batch\_retrieve\_polymorphs}} \\
\begin{description}
    \item[\textbf{Description:}] Retrieve polymorphs for multiple chemical compositions efficiently in batch mode and save it to given directory as json.
    \item[\textbf{Arguments:}]
    \begin{description}
        \item \texttt{compositions} (list[str]): List of chemical compositions to retrieve.
        \item \texttt{max\_energy\_above\_hull} (float): Maximum energy above hull threshold in eV/atom. Defaults to 0.5.
        \item \texttt{max\_per\_composition} (int): Maximum number of polymorphs per composition. Defaults to 10.
        \item \texttt{save\_directory} (str): Directory path for saving individual composition files. Defaults to "polymorph\_data".
    \end{description}
    \item[\textbf{Returns:}] JSON string with batch retrieval results and statistics.

A comprehensive JSON report containing lists of successfully processed and failed compositions, total number of polymorphs retrieved, file paths for each composition, and summary statistics.
         This enables quality control and tracking of the batch processing workflow.
\end{description} \\
\midrule
\textbf{\texttt{cat\_files}} \\
\begin{description}
    \item[\textbf{Description:}] Concatenate and display contents of one or more files
    \item[\textbf{Arguments:}]
    \begin{description}
        \item \texttt{paths} (list[str]): List of file paths to concatenate
        \item \texttt{separator} (str): Separator between file contents
    \end{description}
\end{description} \\
\midrule
\textbf{\texttt{consolidate\_polymorph\_datasets}} \\
\begin{description}
    \item[\textbf{Description:}] Consolidate multiple polymorph JSON files into a single comprehensive dataset.
    \item[\textbf{Arguments:}]
    \begin{description}
        \item \texttt{composition\_files} (dict[str, str]): Dictionary mapping compositions to their JSON file paths.
        \item \texttt{output\_path} (str): Path for the consolidated dataset file. Defaults to "consolidated\_polymorphs.json".
    \end{description}
    \item[\textbf{Returns:}] JSON string with consolidation results and comprehensive statistics.

A JSON-formatted string containing consolidation status, output file path, and detailed statistics including total number of polymorphs, number of compositions successfully included, average polymorphs per composition, and any processing errors.
         This enables quality control and assessment of the consolidation process.
\end{description} \\
\midrule
\textbf{\texttt{copy\_file}} \\
\begin{description}
    \item[\textbf{Description:}] Copy a file from a source to destination
    \item[\textbf{Arguments:}]
    \begin{description}
        \item \texttt{source} (str): Source file path
        \item \texttt{destination} (str): Destination file path
    \end{description}
\end{description} \\
\midrule
\textbf{\texttt{evaluate\_xgboost\_model}} \\
\begin{description}
    \item[\textbf{Description:}] Evaluate trained XGBoost model with comprehensive performance metrics and detailed analysis.
    \item[\textbf{Arguments:}]
    \begin{description}
        \item \texttt{model\_path} (str): Path to the saved XGBoost model file.
        \item \texttt{test\_data\_path} (str): Path to test data CSV file with same structure as training data.
        \item \texttt{target\_column} (str): Name of the target column for evaluation. Defaults to "formation\_energy\_per\_atom".
        \item \texttt{detailed\_analysis} (bool): Whether to include detailed analysis and feature importance. Defaults to True.
    \end{description}
    \item[\textbf{Returns:}] JSON string with comprehensive evaluation metrics and analysis results.

A detailed JSON-formatted string containing evaluation success status, comprehensive performance metrics (MAE, RMSE, R2, MAPE), prediction statistics (min/max/std), error analysis (mean error, error std, max positive/negative errors), and top 10 feature importance rankings when detailed analysis is enabled.
         This provides complete model assessment for validation and comparison purposes.
\end{description} \\
\midrule
\textbf{\texttt{file\_info}} \\
\begin{description}
    \item[\textbf{Description:}] Get information about a file or directory
    \item[\textbf{Arguments:}]
    \begin{description}
        \item \texttt{path} (str): Path to the file or directory
    \end{description}
\end{description} \\
\midrule
\textbf{\texttt{filter\_json\_with\_strategy}} \\
\begin{description}
    \item[\textbf{Description:}] Filter JSON data using custom Python code and save results to a new file.
    \item[\textbf{Arguments:}]
    \begin{description}
        \item \texttt{input\_json\_path} (str): Path to the input JSON file to be filtered.
        \item \texttt{output\_json\_path} (str): Path where filtered JSON data will be saved.
        \item \texttt{custom\_code} (str | None): Python code string defining the filtering logic.
    \end{description}
    \item[\textbf{Returns:}] JSON string with filtering results and comprehensive statistics.

A JSON-formatted string containing filtering status, original and filtered data counts, output file path, and percentage reduction achieved.
         This provides comprehensive information about the filtering operation's effectiveness and enables quality control of the data processing pipeline.
\end{description} \\
\midrule
\textbf{\texttt{get\_bulk\_polymorphs\_data}} \\
\begin{description}
    \item[\textbf{Description:}] Query Materials Project database to find all polymorphs for a given chemical composition.
    \item[\textbf{Arguments:}]
    \begin{description}
        \item \texttt{composition} (str): Chemical composition formula.
    \end{description}
    \item[\textbf{Returns:}] JSON string containing comprehensive polymorph data sorted by stability.

A JSON-formatted string containing a list of dictionaries, each representing a polymorph with properties including Materials Project ID, CIF structure, energy above hull, formation energy per atom, band gap, density, volume, number of sites, space group, and stability information.
         Results are sorted by energy above hull for easy identification of the most stable phases.
\end{description} \\
\midrule
\textbf{\texttt{get\_bulk\_polymorphs\_data\_to\_file}} \\
\begin{description}
    \item[\textbf{Description:}] Query Materials Project for polymorphs and save comprehensive data to a JSON file to give path.
    \item[\textbf{Arguments:}]
    \begin{description}
        \item \texttt{composition} (str): Chemical composition formula.
        \item \texttt{save\_path} (str | None): File path where JSON data will be saved.
    \end{description}
    \item[\textbf{Returns:}] File path where the polymorph data was saved.

Returns the exact file path where the JSON data was successfully written.
         This path can be used by subsequent tools for data loading and processing.
         The file contains comprehensive polymorph data in JSON format, sorted by thermodynamic stability.
\end{description} \\
\midrule
\textbf{\texttt{list\_files}} \\
\begin{description}
    \item[\textbf{Description:}] List files in a directory
    \item[\textbf{Arguments:}]
    \begin{description}
        \item \texttt{path} (str): Path to the directory
        \item \texttt{recursive} (bool): Whether to list files recursively
    \end{description}
\end{description} \\
\midrule
\textbf{\texttt{perform\_cross\_validation}} \\
\begin{description}
    \item[\textbf{Description:}] Perform k-fold cross-validation to assess model stability and generalization performance.
    \item[\textbf{Arguments:}]
    \begin{description}
        \item \texttt{train\_data\_path} (str): Path to training data CSV file.
        \item \texttt{target\_column} (str): Name of the target column for prediction. Defaults to "formation\_energy\_per\_atom".
        \item \texttt{cv\_folds} (int): Number of cross-validation folds. Defaults to 5.
        \item \texttt{hyperparameters} (dict | None): Optional XGBoost hyperparameters for cross-validation.
    \end{description}
    \item[\textbf{Returns:}] JSON string with comprehensive cross-validation results and statistical analysis.

A detailed JSON-formatted string containing cross-validation success status, individual fold scores, statistical summaries (mean, standard deviation), hyperparameters used, and performance stability assessment.
         This enables comprehensive evaluation of model robustness and generalization capability.
\end{description} \\
\midrule
\textbf{\texttt{prepare\_tabular\_dataset}} \\
\begin{description}
    \item[\textbf{Description:}] Prepare tabular dataset for traditional ML models with feature engineering.
    \item[\textbf{Arguments:}]
    \begin{description}
        \item \texttt{polymorphs\_json\_path} (str): Path to consolidated polymorphs JSON file.
        \item \texttt{output\_path} (str): Base path for saving dataset files.
        \item \texttt{target\_property} (str): Property to predict. Defaults to "formation\_energy\_per\_atom".
        \item \texttt{feature\_engineering} (str): Feature engineering strategy. Defaults to "basic".
        \item \texttt{test\_split} (float): Fraction of data for test set. Defaults to 0.2.
        \item \texttt{normalize} (bool): Whether to normalize features. Defaults to True.
    \end{description}
    \item[\textbf{Returns:}] JSON string with comprehensive dataset preparation results and file paths.

A detailed JSON-formatted string containing preparation success status, file paths for training and test data, normalization parameters, dataset statistics, feature information, and metadata.
         This provides complete information about the prepared dataset for subsequent ML workflows.
\end{description} \\
\midrule
\textbf{\texttt{read\_file}} \\
\begin{description}
    \item[\textbf{Description:}] Read the contents of a file into a string
    \item[\textbf{Arguments:}]
    \begin{description}
        \item \texttt{path} (str): Path to the file to read
    \end{description}
\end{description} \\
\midrule
\textbf{\texttt{select\_polymorphs\_with\_strategy}} \\
\begin{description}
    \item[\textbf{Description:}] Select polymorphs using strategic criteria for systematic materials analysis.
    \item[\textbf{Arguments:}]
    \begin{description}
        \item \texttt{polymorphs\_data} (str): JSON string or file path containing polymorph data.
        \item \texttt{selection\_strategy} (str): Strategy for polymorph selection. Defaults to "diverse\_energy".
        \item \texttt{max\_polymorphs} (int): Maximum number of polymorphs to select. Defaults to 5.
        \item \texttt{energy\_threshold} (float): Maximum energy above hull in eV/atom. Defaults to 0.5.
        \item \texttt{is\_path} (bool): Whether polymorphs\_data is a file path. Defaults to False.
    \end{description}
    \item[\textbf{Returns:}] JSON string containing selected polymorphs based on the specified strategy.

A JSON-formatted string containing the selected subset of polymorphs, maintaining the same data structure as the input but with reduced number of entries.
         The selection preserves important characteristics according to the chosen strategy while reducing dataset size for efficient processing.
\end{description} \\
\midrule
\textbf{\texttt{select\_polymorphs\_with\_strategy\_to\_file}} \\
\begin{description}
    \item[\textbf{Description:}] Select polymorphs using strategic criteria and save results to file for persistent storage.
    \item[\textbf{Arguments:}]
    \begin{description}
        \item \texttt{polymorphs\_data} (str): JSON string or file path containing polymorph data.
        \item \texttt{save\_path} (str): File path where selected polymorphs will be saved.
        \item \texttt{selection\_strategy} (str): Strategy for polymorph selection. Defaults to "diverse\_energy".
        \item \texttt{max\_polymorphs} (int): Maximum number of polymorphs to select. Defaults to 5.
        \item \texttt{energy\_threshold} (float): Maximum energy above hull in eV/atom. Defaults to 0.5.
        \item \texttt{is\_path} (bool): Whether polymorphs\_data is a file path. Defaults to False.
    \end{description}
    \item[\textbf{Returns:}] File path where the selected polymorphs were saved.

The complete file path where the selected polymorphs have been successfully saved.
         This path can be used by subsequent tools for loading the selected dataset or for verification that the file was created correctly.
         The file contains the subset of polymorphs selected according to the specified strategy.
\end{description} \\
\midrule
\textbf{\texttt{sort\_and\_get\_first\_from\_json}} \\
\begin{description}
    \item[\textbf{Description:}] Sort JSON data by specified key and return the first element's specified value.
    \item[\textbf{Arguments:}]
    \begin{description}
        \item \texttt{polymorph\_data\_json} (str): JSON string containing the data to be sorted.
        \item \texttt{sort\_key} (str): Property name to sort the data by.
        \item \texttt{return\_key} (str): Property name to return from the first element after sorting.
    \end{description}
    \item[\textbf{Returns:}] Value of the specified return\_key from the first element after sorting.

The value corresponding to the return\_key from the material that has the smallest value for the sort\_key. This could be a string (like material\_id or CIF), a number (like energy or band gap), or any other data type stored in the JSON. The returned value represents the optimal material according to the specified sorting criterion.
\end{description} \\
\midrule
\textbf{\texttt{train\_xgboost\_model}} \\
\begin{description}
    \item[\textbf{Description:}] Train XGBoost regression model for property prediction with evaluation.
    \item[\textbf{Arguments:}]
    \begin{description}
        \item \texttt{train\_data\_path} (str): Path to training data CSV file.
        \item \texttt{test\_data\_path} (str): Path to test data CSV file.
        \item \texttt{model\_save\_path} (str): Path to save the trained model file in .pkl format
        \item \texttt{target\_column} (str): Name of the target column for prediction. Defaults to "formation\_energy\_per\_atom".
        \item \texttt{hyperparameters} (dict | None): Optional dictionary of XGBoost hyperparameters.
    \end{description}
    \item[\textbf{Returns:}] JSON string with comprehensive training results and model performance metrics.

A detailed JSON-formatted string containing training success status, model performance metrics (MAE, RMSE, R2), feature importance rankings, hyperparameters used, dataset information, and file paths for saved model and results.
         This enables comprehensive model evaluation and comparison.
\end{description} \\
\midrule
\textbf{\texttt{write\_file}} \\
\begin{description}
    \item[\textbf{Description:}] Write content to a file
    \item[\textbf{Arguments:}]
    \begin{description}
        \item \texttt{path} (str): Path to the file to write
        \item \texttt{content} (str): Content to write to the file
    \end{description}
\end{description} \\
\end{longtable}
